% 阋神星（综述）
% license CCBYSA3
% type Wiki

本文根据 CC-BY-SA 协议转载翻译自维基百科\href{https://en.wikipedia.org/wiki/Eris_(dwarf_planet)}{相关文章}。

\begin{figure}[ht]
\centering
\includegraphics[width=6cm]{./figures/88a62f7999189f0d.png}
\caption{Hubble Space Telescope 于 2006 年 8 月拍摄的 Eris 和 Dysnomia 的低分辨率图像} \label{fig_Eris_1}
\end{figure}
Eris（小行星编号：136199 Eris）是太阳系中已知质量最大、体积第二大的矮行星\(^{[23]}\)。它是一颗处于散射盘中的跨海王星天体（TNO），具有高偏心率轨道。Eris 于 2005 年 1 月由 Mike Brown 领导的 Palomar Observatory 团队发现，并于同年稍后得到确认。该天体于 2006 年 9 月以希腊—罗马神话中的纷争与不和女神命名。Eris 是围绕太阳运行的已知天体中质量第九大者，在整个太阳系中（包括卫星）质量排名第十六。它也是太阳系中已知尚未被航天器造访的最大天体。Eris 的直径测量为 2,326 ± 12 km（1,445 ± 7 mi）\(^{[13]}\)；其质量为 Earth 的 0.28\%，比 Pluto 大 27\%\(^{[24][25]}\)，尽管 Pluto 在体积上略大一些\(^{[26]}\)。Eris 与 Pluto 的表面积均可与 Russia 或 South America 相当。

Eris 已知有一颗大型卫星，名为 Dysnomia。2016 年 2 月，Eris 与太阳的距离为 96.3 AU（144.1 亿 km；89.5 亿 mi）\(^{[21]}\)，超过 Neptune 或 Pluto 距离的三倍。除长周期彗星外，直到 2018 年发现 2018 AG37 和 2018 VG18 之前，Eris 与 Dysnomia 一直是太阳系中已知距离最遥远的天然天体\(^{[21]}\)。

由于 Eris 起初看起来比 Pluto 更大，NASA 最初将其称为太阳系的第十颗行星。这一认识及未来可能发现与之大小相近的其他天体的前景，促使国际天文学联合会（IAU）首次定义“行星”一词。根据 IAU 于 2006 年 8 月 24 日批准的定义，Eris、Pluto 和 Ceres 均为“矮行星”\(^{[27]}\)，使太阳系已知行星的数量减少为八颗，与 1930 年 Pluto 被发现之前相同。2010 年一次 Eris 对恒星的掩食观测表明，它略小于 Pluto\(^{[28][29]}\)，而 New Horizons 于 2015 年 7 月测得 Pluto 的平均直径为 2,377 ± 4 km（1,477 ± 2 mi）\(^{[30][31]}\)。
\subsection{发现}  
Eris 由 Mike Brown、Chad Trujillo 和 David Rabinowitz\(^{[2]}\) 于 2005 年 1 月 5 日发现，所使用的图像拍摄于 2003 年 10 月 21 日\(^{[32]}\)。该发现于 2005 年 7 月 29 日公布，与 Makemake 的公布在同一天，且比 Haumea 晚两天\(^{[33]}\)，部分原因与后来有关 Haumea 的争议事件有关。该搜寻团队多年来一直在系统性地扫描太阳系外层的大型天体，并参与发现了数颗其他大型 TNO，包括矮行星 Quaoar、Orcus 和 Sedna\(^{[34]}\)。

该团队于 2003 年 10 月 21 日使用位于 California 的 Palomar Observatory 的 1.2 m Samuel Oschin Schmidt 望远镜进行了常规观测，但当时未能识别出 Eris 的图像，因为它在天空中的运动非常缓慢：团队的自动图像搜索软件会排除所有运动速度小于每小时 1.5 角秒的天体，以减少返回的误报数量\(^{[32]}\)。2003 年 Sedna 被发现时，其运动速度为每小时 1.75 角秒，因此团队以更低的角运动阈值重新分析了旧数据，并人工筛选了此前被排除的图像。2005 年 1 月，重新分析揭示了 Eris 相对于背景恒星的缓慢轨道运动\(^{[32]}\)。

随后进行了后续观测，以便初步确定 Eris 的轨道，从而估计该天体的距离\(^{[32]}\)。团队原计划在进一步观测和计算完成之前，暂缓公布明亮天体 Eris 和 Makemake 的发现，但在团队一直追踪的另一颗大型 TNO——Haumea——于 7 月 27 日被一个西班牙团队有争议地公布后，两个发现均于 7 月 29 日宣布\(^{[2]}\)。

Eris 的回溯预发现图像已被追溯至 1954 年 9 月 3 日\(^{[6]}\)。

2005 年 10 月公布的进一步观测表明，Eris 有一颗卫星，后来被命名为 Dysnomia。对 Dysnomia 轨道的观测使科学家能够确定 Eris 的质量，并在 2007 年 6 月计算为 \((1.66\pm0.02)\times10^{22}\) kg\(^{[24]}\)，比 Pluto 大 27\%±2\%。
\subsection{名称}  
Eris 的名称来源于希腊女神 Eris（希腊语 Ἔρις），她是纷争与不和的拟人化形象\(^{[35]}\)。该名称由 Caltech 团队于 2006 年 9 月 6 日提出，并于 2006 年 9 月 13 日正式获准\(^{[36]}\)，这一命名在此之前经历了异常漫长的阶段，在这期间该天体一直以临时编号 2003 UB313 被称呼，此编号是 IAU 根据小行星命名规则自动授予的。

名称 Eris 存在两种竞争性读法，分别使用“长元音”或“短元音”，类似单词 era 的两种不同读法\(^{[4]}\)。在英语中较为常见的形式（例如 Brown 及其学生使用的读法）是 /ˈɛrɪs/，为双音节、短元音读法\(^{[3]}\)。然而，该女神的传统英语读法为 /ˈɪərɪs/，带有长元音\(^{[5]}\)。

该名称在希腊语和拉丁语中的斜格词干为 Erid-\(^{[38]}\)，如意大利语 Eride 和俄语 Эрида（Erida）所示，因此英语形容词形式为 Eridian /ɛˈrɪdiən/\(^{[9][10]}\)。
\subsubsection{Xena}  
由于当时不确定该天体被归类为行星还是小行星（因为这两类天体适用不同的命名程序）\(^{[39]}\)，因此必须等到 2006 年 8 月 24 日 IAU 做出决定后才能正式命名\(^{[40]}\)。在此期间，该天体在公众中一度被称为 Xena。“Xena” 是发现团队内部使用的非正式昵称，灵感来自电视剧《战士公主：西娜》（Xena: Warrior Princess）的主角。据报道，发现团队一直保留“Xena”这一昵称，用于命名第一颗被发现且体积大于 Pluto 的天体。据 Brown 所说：

“我们选择它是因为它以 X 开头（即行星‘X’），听起来很神话……同时我们也希望更多女性神祇的名字能被使用（如 Sedna）。此外，当时该电视剧仍在播出，这说明我们已经搜索了相当长时间！”\(^{[41]}\)

Brown 在一次采访中表示命名过程曾一度陷入停滞：

“一位记者（Ken Chang）\(^{[42]}\) 给我打电话，他当时在《纽约时报》工作，碰巧是我大学同学……他问我：‘你们提议的名字是什么？’我说：‘这个我不能告诉你。’然后他又问：‘那你们私下互相交流时叫它什么？’……据我所知，那是我唯一一次在媒体上说出这个名字，结果传播开来，我对此只有一点点愧疚；其实我还挺喜欢这个名字。”\(^{[43]}\)
\subsubsection{选择正式名称}
\begin{figure}[ht]
\centering
\includegraphics[width=6cm]{./figures/2efad5f955645faf.png}
\caption{} \label{fig_Eris_3}
\end{figure}















\subsection{另见}
\begin{itemize}
\item 天体命名
\item 清除邻近的小天体
\item 国际天文学联合会（IAU）
\item 假设的海王星外行星
\end{itemize}
\subsection{参考文献}
\begin{enumerate}
\item Staff. Discovery Circumstances: Numbered Minor Planets. IAU: Minor Planet Center. 2007-05-01 [2007-05-05]. （原始内容存档于2016-05-04）.
\item Brown, Mike. The discovery of 2003 UB313 Eris, the largest known dwarf planet. 2006 [2007-05-03]. （原始内容存档于2014-09-10）.
\item JPL Small-Body Database Browser: 136199 Eris (2003 UB313). 2009-11-20 [2009-01-21]. （原始内容存档于2019-01-09）.
\item List Of Centaurs and Scattered-Disk Objects. Minor Planet Center. [2008-09-10]. （原始内容存档于2011-07-25）.
\item Buie, Marc. Orbit Fit and Astrometric record for 136199. Deep Ecliptic Survey. 2007-11-06 [2007-12-08]. （原始内容存档于2014-09-26）.
\item Asteroid Observing Services （页面存档备份，存于互联网档案馆）
\item Sicardy, B.; Ortiz, J. L.; Assafin, M.; Jehin, E.; Maury, A.; Lellouch, E.; Gil-Hutton, R.; Braga-Ribas, F.; Colas, F.; Widemann. Size, density, albedo and atmosphere limit of dwarf planet Eris from a stellar occultation (PDF). European Planetary Science Congress Abstracts. 2011, 6: 137 [September 14, 2011]. Bibcode:2011epsc.conf..137S. （原始内容存档 (PDF)于October 18, 2011）.
\item Beatty, Kelly. Former 'tenth planet' may be smaller than Pluto. NewScientist.com. Sky and Telescope. November 2010 [October 17, 2011]. （原始内容存档于February 23, 2012）.
\item 掩星觀測顯示鬩神星和冥王星大小幾乎相等. PanSci 泛科学. 2011-11-03 [2017-03-09]. （原始内容存档于2014-03-25） （中文（台湾））.
\item Michael E. Brown and Emily L. Schaller. The Mass of Dwarf Planet Eris (abstract page). Science. 2007, 316 (5831): 1585 [2011-02-01]. PMID 17569855. doi:10.1126/science.1139415. （原始内容存档于2010-06-26）.
\item Kelly Beatty. Eris Gets Dwarfed (Is Pluto Bigger?). Sky & Telescope (News Blog). 2010-11-07 [2010-11-07]. （原始内容存档于2011-08-30）.
\item Snodgrass, Carry, Dumas, Hainaut. Characterisation of candidate members of (136108) Haumea's family (abstract). The Astrophysical Journal. 2009-12-16. arXiv:0912.3171.
\item AstDys (136199) Eris Ephemerides. Department of Mathematics, University of Pisa, Italy. [2009-03-16]. （原始内容存档于2011-06-04）.
\item F. Bertoldi, W. Altenhoff, A. Weiss, K.M. Menten, C. Thum. The trans-neptunian object UB313 is larger than Pluto. Nature. 2006-02, 439 (7076): 563–564 [2018-04-02]. ISSN 1476-4687. doi:10.1038/nature04494. （原始内容存档于2019-07-13） （英语）.
\item “阋”，《说文解字》卷三：“恒讼也。《诗》云：‘兄弟阋于墙。’从斗从儿。儿，善讼者也。许激切。”繁：“鬩”，简：“阋”，拼音：xì，注音：ㄒㄧˋ，粤拼：jik1
\item JPL/NASA. What is a Dwarf Planet?. Jet Propulsion Laboratory. 2015-04-22 [2022-01-19]. （原始内容存档于2021-12-08）.
\item How Big Is Pluto? New Horizons Settles Decades-Long Debate. www.nasa.gov. 2015 [2015-07-14]. （原始内容存档于2017-07-01）.
\item Mike Brown. The discovery of 2003 UB313 Eris, the 10th planet largest known dwarf planet. Caltech. 2006 [2010-01-05]. （原始内容存档于2012-05-20）.
\item The IAU draft definition of "planet" and "plutons" (新闻稿). IAU. 2006-08-16 [2006-08-16]. （原始内容存档于2006-08-20）.
\item Mike Brown. The shadowy hand of Eris. Mike Brown's Planets. 2010 [2010-11-07]. （原始内容存档于2010-11-11）.
\item Mike Brown. How big is Pluto, anyway?. Mike Brown's Planets. 2010-11-22 [2010-11-23]. （原始内容存档于2015-10-31）. (Franck Marchis on 2010-11-08)
\item Brown, Mike. How big is Pluto, anyway?. Mike Brown's Planets. 2010-11-22 [2010-11-23]. （原始内容存档于2015-10-31）. (Franck Marchis on 2010-11-08)
\item Thomas H. Maugh II and John Johnson Jr. His Stellar Discovery Is Eclipsed. Los Angeles Times. 2005-10-16 [2008-07-14]. （原始内容存档于2017-02-21）.
\item KLC. 矮行星Eris中譯名正式命名為「鬩神星」. tamweb.tam.gov.tw. 台北天文馆. 2007-07-10 [2017-03-09]. （原始内容存档于2013-07-29） （中文（台湾））.
\item 方舟子. 方舟子：欺世盗名的八卦宇宙论. 人民网. 2005-08-01 [2013-03-30]. （原始内容存档于2013-05-09） （中文（中国大陆））.
\item Gemini Observatory Shows That"10th Planet" Has a Pluto-Like Surface. Gemini Observatory. 2005 [2007-05-03]. （原始内容存档于2007-03-11）.
\item M. E. Brown, C. A.Trujillo, D. L. Rabinowitz. Discovery of a Planetary-sized Object in the Scattered Kuiper Belt. The Astrophysical Journal (abstract page). 2005, 635 (1): L97–L100. Bibcode:2005ApJ...635L..97B. arXiv:astro-ph/0508633 可免费查阅. doi:10.1086/499336.
\item J. Licandro, W. M. Grundy, N. Pinilla-Alonso, P. Leisy. Visible spectroscopy of 2003 UB313: evidence for N2 ice on the surface of the largest TNO (PDF). Astronomy and Astrophysics. 2006, 458 (1): L5–L8 [2011-05-03]. Bibcode:2006A&A...458L...5L. arXiv:astro-ph/0608044 可免费查阅. doi:10.1051/0004-6361:20066028. （原始内容存档 (PDF)于2008-11-21）.
\item M. E. Brown, M. A. van Dam, A. H. Bouchez, D. LeMignant, C. A. Trujillo, R. Campbell, J. Chin, Conrad A., S. Hartman, E. Johansson, R. Lafon, D. L. Rabinowitz, P. Stomski, D. Summers, P. L. Wizinowich. Satellites of the largest Kuiper belt objects. The Astrophysical Journal (abstract page). 2006, 639 (1): L43–L46. Bibcode:2006ApJ...639L..43B. arXiv:astro-ph/0510029 可免费查阅. doi:10.1086/501524.
\end{enumerate}
\subsection{外部链接}
\begin{itemize}
\item Michael Brown 关于 Eris 和 Dysnomia 的网页  
\item MPC 关于 (136199) Eris 的数据库条目  
\item NASA JPL 的 3D 轨道可视化  
\item 2003 UB313 轨道模拟  
\item Keck 天文台关于 Dysnomia 发现的页面  
\item AstDyS-2（小行星动力学网站）中的 Eris（矮行星）  
\item 历表 · 观测预测 · 轨道信息 · 本征要素 · 观测信息  
\item JPL 小天体数据库中的 Eris（可在 Wikidata 编辑）  
\item 近距离接近 · 发现 · 历表 · 轨道查看器 · 轨道参数 · 物理参数
\end{itemize}












